\documentclass[12pt]{article}
\usepackage[utf8]{inputenc}
\usepackage{amsmath, amssymb} % Math packages
\usepackage{geometry}
\usepackage{hyperref} % Clickable links
\usepackage{graphicx} % Images

\geometry{a4paper, margin=1in}

\title{reti di calcolatori }
\author{Matteo}
\date{\today}

\begin{document}


\maketitle
\section{intro}

\textbf{reti di calcolatori}: un insieme di dispositivi di calcolo, autonomi o interconnessi.

Una rete di calcolatori è un insieme di dispositivi autonomi, cioè in grado di eseguire e svolgere autonomamente i compiti
programmati di calcolo e di comunicazione, interconnessi tra loro da supporti fisici alla trasmissione di segnali.
i televisori e le reti telefoniche ad esempio non sono reti di calcolatori dato che non sono autonomi

Le reti di calcolatori si possono classificare in base alla dimensione geografica:
\begin{itemize}
\item Reti personali (Personal Area Network), PAN
 Connessione tra dispositivi in una stanza, o di una scrivania
\item Reti locali (Local Area Network), LAN
 Connessione dispositivi di un ufficio, un laboratorio, un edificio, un campus
\item Reti metropolitane (Metropolitan Area Network), MAN
Connessione di aree urbane, reti civiche
\item Reti geografiche (Wide Area Network), WAN
 Connessione di aree molto ampie, reti nazionali e internazionali
\item Internet: è la rete globale composta dall’unione di reti di vari tipi, connesse tra loro e conformi a un
determinato insieme di regole di comunicazione comuni (i protocolli di Internet).
\end{itemize}

lo sviluppo delle reti di calcolatori che ha permesso la nascita di internet non sarebbe stata possibile senza una 
distribuzione dei costi delle infrastrutture, che e avvenuta in modo organico per rispondere a esigenze crescenti degli utenti.

due grandi aspetti di interesse degli utenti sono \textbf{capacita di trasmissione}, misurati in byte(kilo,mega giga tera) e \textbf{ritardo del collegamento}(ad esempio jitter,la variazione della latenza intorno al val medio)

il ritado del collegamento indica il tempo effettivoo per transitare dal mittente al destinatario finale sulla rete, il ritardo puo essere dovuto a distanza fisica o a tempo impiegato a gestire i protocolli di rete.

la connessione di un calcolatore a una rete richiede un minimo di componenti:
\begin{itemize}
\item dispositivi o schede di rete: sono dispositivi hardware per codificare, trasmettere, ricevere e decodificare
dati dal calcolatore alla rete, e viceversa
, amministrati da componenti software del sistema operativo
\item mezzo di trasmissione: supporto fisico alla propagazione e trasmissione di segnali: es. cavi elettrici,
fibre ottiche. Esso realizza l’infrastruttura fisica di rete.
\item connettore di rete: interfaccia o connettore standard per il collegamento del dispositivo di rete al mezzo
di trasmissione della rete
\item protocolli di rete: insieme di regole implementate sotto forma di software del calcolatore, univocamente
definite per garantire compatibilità, e corretta gestione della comunicazione
\end{itemize}

\textbf{collegamento di rete}: mezzo di trasmissione fisico condiviso tra calcolatori in rete

\textbf{infrastruttura di rete}: struttura di connessione dei collegamenti di rete, puo essere:
\begin{itemize}
\item punto a punto(2 calcolatori)
\item connessioni completamente connesse o cammini ridondanti(ogni calcolatore econnesso a ogni altro)
\item connessioni parzialmente connesse(ogni calcolatore e connesso a un altro solo, ma la rete li prende tutti)
\item coonessioni a partizioni di rete (uno o piu gruppi e separato dagli altri)
\end{itemize}

nelle connessioni completamente connesse quel numero di connessioni e anche eccessivo, ma e mollto resistente ai rischi.
nelle connessioni parzialmente connesse invece ce il miglior rapporto costo efficacio ma un solo guasto comprometterebbe l'intera rete.


Le infrastrutture di rete possono assumere vari schemi di connessione dei dispositivi \textbf{topologia}

varii tipi di topologie
\begin{itemize}
\item a cerchio, ogni calc  e connesso solo al precedente e successivo, i dati partono in 1 senso e tornano nell'altro
\item a stella, ogni elemento periferico e connesso bidirezionalmente a un unico calcolatore centrale. l'elemento centrale e detto "single point failure"
\item a bus, ogni elemento ha una connessione con il bus centrale
\item ad albero, abbiamo elementi padri che si connettono ad elementi figli in uno schema ad albero.
\end{itemize}

nel bus ce il problema di coordinare chi puo usare il bus in che momento per evitare sovrapposizioni di seganle.

man mano che le infrastrutture si fanno piu complesse questi schemi diventano sempre meno applicabili.

per capirsi tra calcolatori con codifica non return zero(la piu semplice) era necessario essere perfettamnete simcronizzati coi clock cosi fu inventata la codifica manchester

queste codifiche leggono ilsegnale elettrico come 1 e l’assenza come 0 

la manchester riallinea se i clock si desincronizzano

vecchi computer usavano la manchester ora non e piu un problema


i \textbf{mezzi di trasmissione} possono essere principalmente di 3 tipi:
\begin{itemize}
\item cavetti o fili metallici, come i doppini intrecciati
\item fibra ottica, trasmissione segnali luminosi vincolati entro fibre di vetro
\item connessioni wireless,  radiazioni elettromagnetiche nello spazio, non sempre affidabili
\end{itemize}


\textbf{dispositivo o scheda di rete}:
dotato di interfaccia di collegamento al calcolatore e connettore di rete

codifica e decodifica dati in ricezione e trasmissione

esistono schede di rete diverse per mezzi fisici diversi(cablati,fibra,wifi)

Prende il nome dai protocolli o dallo standard utilizzato per la trasmissione

ognuno ha un codice identificativo unico \textbf{indirizzo MAC}(medium access control)


canale di comunicazione della rete: una visione astratta del mezzo di trasmissione

puo essere da punto a punto

o un canale ad accesso multiplo(broadcast):
tutti i dispositivi ricevono sullo stesso canale. e richiesto arbitragiio per evitare collisioni.
e richiesto anche indirizzamento Indirizzi univoci per mittente e destinatario (o destinatari) del dato trasmesso e indirizzo MAC della scheda di rete.


Una serie di canali in cascata può realizzare un servizio di trasferimento dell’informazione tra due dispositivi
anche molto distanti, secondo due modalità principali: \textbf{commutazione di circuito} (o circuito riservato) e
\textbf{commutazione di pacchetto.}
Commutazione di circuito: ad esempio, linea telefonica Viene riservato un circuito di canali di comunicazione punto a punto su ogni connessione lungo
tutto il cammino dal mittente al destinatario

Il circuito riservato ha un ritardo di comunicazione basso

Si paga il tempo di connessione, anche se non si scambiano dati sul canale

Basso utilizzo delle risorse di rete (canale sprecato quando non si trasmette)

dopo averlo fatto la prima volta non ce bisogno di dichiarare periodicamente mittente e destinatario
\\
commutazione di pacchetto: molto usata su internet e canali broadcast.

i dati vengono divisi in pacchetti indipendenti con sopra indirizzi di destinatario e mittente

poi ogni pacchetto viene spedito separatamente

diversi flussi di pacchetti condividono il canale, utilizzandolo al massimo. i nodi ricevitori devono di volta in volta
farsi carico di verificare se il pacchetto sia giunto a destinazione o, in caso contrario, possono provvedere all’inoltro del
pacchetto ricevuto sul successivo canale broadcast.

richiede quindi meno risorse ma ha un ritardo della rete di comunicazione.

inoltre  si paga in base alla quantita di dati trasmessi.

\subsection{funzionamento della commutazione a pacchetti}

nella commutazione a pacchetti a ogni nodo il pacchetto viene inoltrato verso il destinatario finale e possibile la perdita o arrivo disordinato dei pacchetti.

esistono sia: \textbf{servizi orientati alla connessione} che fanno in modo i pacchetti arrivino sempre in ordine e garantiscono la ritrasmissione di pacchetti perduti
che \textbf{servizi non orientati alla connessione} in cui invece i pacchetti possono seguire strade diverse e arrivare tardi o non arrivare mai.

definizione \textbf{protocollo}:Regole e procedure (semantiche) di gestione dei processi di comunicazione
 Regole e formati (sintattici) per definire scambio di messaggi non ambigui
 permettono compatibilità dei dispositivi e dei sistemi conformi allo stesso standard

 definizione \textbf{architettura dei protocolli di rete}:
 ogni livello risolve uno dei problemi della comunicazione.
 i livelli superiori effettuano richieste di servizio ai livelli inferiori.
i livelli inferiori forniscono servizi ai livelli superiori.
Le richieste e i servizi realizzano l’interfaccia del protocollo verso altri livelli.
ogni livello comunica solo con il livello superiore e inferiore.

c'è un modello standard per definire l'architettura dei protocolli di rete, si chiama \textbf{Standard ISO/OSI RM} (Open System Interconnection Reference Model):

\begin{enumerate}
\item[7.] livello applicazione, primitiv dati applicazione
\item[6.] livello presentazione, traduzione formati dati utente
\item[5.] livello sessione, stato del collegamento
\item[4.] livello trasporto, connessione affidabile e controllo congestione
\item[3.] livello rete, frammentazione, indirizzamento, instradamento pacchetti
\item[2.] livello mac/llc, accesso al medium controllo errore
\item[1.] livello fisico, codifica trasmisssione ricezione
\end{enumerate}

Il livello applicazione fornisce i servizi di trasmissione dati alle applicazioni in esecuzione e l permette di comunicare.

il livello presentazione risolve eventuali eterogeneità dei dati tra i nodi della rete.

il livello sessione mantiene lo stato del collegamento. 

il livello trasporto fornisce un servizio di connessione affidabile e controllo della
congestione. 

il livello rete si occupa di frammentazione dei dati, indirizzamento e instradamento dei
pacchetti. 

il livello LLC/MAC si occupa di garantire una comunicazione affidabile e l’arbitraggio di accesso
al mezzo trasmissivo. 

il livello fisico si occupa di codificare i dati e trasmetterli sul medium di
trasmissione.


L’architettura dei livelli di protocolli di Internet utilizza solo 5 dei 7 livelli ISO/OSI RM, il livello presentazione e il livello sessione sono saltati.

livello sessione si preoccupa se il pacchetto e arrivato corretttamnente e gestisce il ritmo di invio dei pacchetti cioe se inviarli a mitraglia o aspettare tra uno e l’altro

livello sessione e presentazione sono oggi obsoleti: livello sessione mntiene stato collegamneton tra mittente e destinatario, se una sessione si interrompe ti fa ripartire da dove eri rimasto, senza il livello trasporto riparte da capo. oggi le connessioni sono talmente veloci che spesso non serve

\underline{in trasmissione}, ogni livello riceve dati dai livelli superiori e li inserisce (incapsula) in “buste” con dati
aggiuntivi, utili ad istruire il corrispondente livello del dispositivo ricevente

\underline{in ricezione} ogni livello verifica i dati della busta, agisce di conseguenza, e passa il contenuto della busta
ai livelli superiori

Il livello trasporto imbusta i dati aggiungendo informazioni utili all’ordinamento e controllo della
velocità di invio delle buste

Il livello rete frammenta i dati in pacchetti, decide il cammino sul quale inviare il pacchetto a
seconda dell’indirizzo del destinatario

Il livello MAC/LLC esegue la consegna finale dei dati a dispositivi di una rete locale.

la parte iniziale della busta e detta \underline{header}, quella posteriore \underline{trailer}.


Il livello uno (fisico) si occupa delle regole per codificare e trasmettere i dati come segnali sul mezzo trasmissivo a disposizione.

il livello mac deve determinare a quali dispositivi i dati sono destinati qui la rete assume i connotati di una lanciano

Il livello tre (rete) affronta il problema dell’integrazione generalizzata di reti, fino a formare reti di reti,
organizzate in gerarchie di reti e sottoreti i router si occupano della gestione dell'instradamento e smistano i pacchetti tra le reti, qui il servizio
di consegna dei pacchetti è \underline{connectionless}(non orientato alla connessione)

Il livello quattro (trasporto) prevede la gestione delle regole che permettano di ripristinare un servizio orientato alla connessione,A questo livello esiste una rete di reti affidabile di tipo orientato alla connessione. Internet è una rete di questo tipo, basata su
due protocolli in particolare, TCP e IP.

Il livello sette (applicazione) mette semplicemente a disposizione il servizio di comunicazione della rete di reti alle applicazioni
dell’utente.

\subsection{livello fisico, codifica dei dati digitali}

codifica e decodifica bit a segnali analogici e viceversa, il valore massimo di bit/secondo trasmessi
è determinata dal tempo necessario alla codifica del valore di ogni bit, poi i segnalli elettromagnetici viaggiano alla velocita della luce.

\subsection{livello mac}

segmento di rete locale: un mezzo di trasmissione condiviso con canale ad accesso multiplo, tutte le schede di rete ricevono le trasmissioni fatte sul mezzo

\underline{il livello mac/llc} si occupa dell'arbitraggio di chi usa il mezzo trasmissivo quando, di indicare a quale scheda di rete e destinato il frame trasmesso,
e se il frame è destinato a lui aa riceverlo e a passarlo al livello superiore.

il livello mac si occupa anche di assicurarsi che il frame(pacchetto) sia stato ricevuto correttamente, se il destinatario ha ricevuto il pachhetto deve inviare un frame di conferma detto aknowledgement,
se il mittente non riceve nulla, riinvia il frame finche non riceve la conferma del successo.


Tre esempi di protocolli di livello due, che possono essere usati in alternativa tra loro per le schede di rete di
segmenti di rete locale, sono:
\begin{itemize}
\item ethernet, La scheda di rete ascolta il canale e trasmette solo se nessuno sta già trasmettendo
, Se viene rilevata una collisione la trasmissione è interrotta per riprovare più tardi
\item IEEE 802.11(wifi) La scheda di rete ascolta il canale e trasmette solo se nessuno sta già trasmettendo
, si cerca di prevenire le collisioni dilazionando le trasmissioni nel tempo
\item Token ring
usata se i dispositivi sono collocati su topologia ad anello cablato  Esiste un frame detto token che viene passato come un “testimone” tra i dispositivi
 Solo chi detiene il token ha diritto di trasmettere, poi deve passare il token
\end{itemize}

ci sono diversi dispositivi utili a comporre reti locali.
\begin{itemize}
\item ripetitore(repeater): da solo il segnale fisico degrada percorrendo distanza fisica, il ripetitore rigenera e amplifica il segnale
\item Hub, punto di contatto delle connessioni a stella, semplicemente un ripetitore con tante entrate e uscite
\item bridge(ponte)  un dispositivo che fa anche da traduttore a livello mac, permette di collegare dispositivi con protocolli mac diversi, I bridge fanno quindi da traduttori dei frame nei formati richiesti dal livello MAC di ogni segmento connesso.
\item switch(commutatore) come il bridge ma permette di collegare molti segmenti Ha capacità di filtrare e inviare i frame sul segmento giusto, leggendo l’indirizzo MAC del
destinatario del frame
\end{itemize}

ripetitore : hub = bridge : switch





































\section{sicurezza}

abbiamo 4 proprieta per la sicurezza:
\begin{itemize}
\item confidenzialità: abbiamo un sender che cripta i messaggi e un receiver che li decodifica, solo il receiver deve capire
\item autenticazione: il sender e il receiver devono confermare l'identità reciproca.
\item integrita: bisogna controllare che il messaggio non sia modificato
\item disponibilita: il servizio deve essere disponibile agli utenti
\end{itemize}


un intruso puo intercettare i messaggi,leggerli,  modificarli, mandarne di suoi, rubare l'identità o intasare servizi.

noi vogliamo garantire sia la cifrazione del messaggio, che l'identita dell'interlocutore che l'integrità del messaggio.

si possono usare socket cifrati a livello trasporto, se no si puo cifrare anche a livello rete.

cufratura e la prima risorsa di sicurezza che viene usata

la prima arma che abbiamo per difenderci è la cifratura

\subsection{crittografia}

esistono percio schemi di codifica per crittare i messaggi.
sender e receiver hanno entrambi una chiave di decodifica per decifrare i messaggi e si accordano per avere la stessa.

gli intrusi non devono ottenere la chiave, percio esistono diversi standard di codifica per garantire la sicurezza.

gli intrusi hanno solo il testo cifrato con cui lavorare, e ci sono due approcci di attacco:

    \textbf{brute force}: si provano tutte le chiavi una alla volta

    \textbf{analisi statistica}

    esistono due tipi di chiavi:
    \begin{itemize}
    \item simmetriche: i due interlocutori condividono la stessa chiave
    \item asimmetriche
    \end{itemize}

DES: data encryption standard. chiavi simmetriche, codifica della chiave inizialmente a 56 bit, poi viene inventato il 3DES, cioe criptare 3 volte con tre chiavi. perche ers diventato troppo facile fare "brute force" fino ad ottenere la chiave.

AES: advanced encryption standard. chiavi simmetriche. processa i dati in blocchi da 128 bit. le chiavi possono essere da 128, 192 o 256 bit. molto piu efficace del DES.

\textbf{cifratura a chiave pubblica}: ha chiave non simmetrica, ognuno ha 2 numeri primi comme chiavi, una solo e privata(segreta) e unica per ognuno, l'altra e pubblica e uguale tra i due. per decifrare devi fare prodotto di due chiavi e messaggio.

debolezza nel fatto che in teoria se il messaggio è troppo ridotto diventa facile decrittarlo.

la chiave privata è fortissima, la debolezza e il messaggio.

quindi esistono i bit di padding, cioe bit senza significato che vengono aggiunti sempre per rendere il messaggio almeno 256 bit, cosi ci sono $2^256$ combinazioni.

la reale forza dell'algoritmo rguarda la fattorizzazione dei numeri primi di grandi dimensioni.

i numeri primi sono divisibili solo per se stessi, crittando usando l valore del prodotto di 2 numeri primi, per trovarlo devi per forza conoscere quei 2 numeri.
quindi fare brute forcing e estremamente difficile.

\subsection{autenticazione}

bisogna rendere sicura annche l'autenticazione perchè altrimenti trudy puo fingersi qualcun altro e ottenre altre informazini.

\textbf{playback attack}: l'attaccante si finge qualcun'altro intercettando A e mandando a B il messaggio intercettato.

un modo per eseguire l'autenticazione prima di cominciare la vera e propria comunicazione e che B mandi ad A un nonce(numero a uso singolo) e che B lo cifri con la sua chiave segreta. questo si chiama AP 4.0 e vale solo con chiavi simmetriche

AP5.0 vale anche con chiavi asimmetriche

\subsection{firme digitali}

una semplice forma di firma e semplicemente crittare con chiave privata

questo pero crea messaggi piuttosto lunghi, quindi viene applicata la funzione hash al messaggio per craere un "messagge digest" di lunghezza fissa e poi lo cripta con la chiave

\textbf{funzione di hashing}: una funzione  e di hashing quando per ogni elemento in un insieme, la destinazione in un altro insieme è casuale.

un paio di funzioni di hash piu usate:
MD5(or RFC 1321): 128 bit di bit digest

SHA-I: produce 160(numero molto inusuale) bit di digest

a questo punto un intruso puo ancora usare una sua chiave pubblica ma dire che e di qualcun altro

entrano quindi in campo le \textbf{certification authority}: enti esterni che associano certe chiavi pubbliche a un individuo, certificandolo.


per quanto riguarda il \textbf{secure e-mail}  si crea una nuova chiave simmetrica, poi sender manda due messaggi, uno il messaggio crittato con chiave simmetrica, la'altro, la chiave simmetrica stessa crittata con la chiave pubblica del receiver. poi manda anche la propria firma digitale come proa di autenticazione.

\subsection{SSL}

SSl o secure socket layer è un protocollo di sicurezza ampiamente usato.

provvede confidenzialità, integrità e autenticazione

è a meta tra il livello applicazione e livello trasporto.

è  praticamente una connessione criptata tra client e server, infatti il SSL protegge le socket che sono i punti da cui i dati escono ed entrano.

SSL è divisa in 4 parti
\begin{enumerate}
    \item handshake sender e receiver usano i loro certificati e chiavi private per autenticarsi e scambiarsi il \underline{shared secret}
    \item key derivation: sender e receiver usano il shared secret per derivare un gruppo di chiavi
    \item data transfer: i dati da trasferire sono divisi in serie di records
    \item connection closure: messaggio speciale per chiudere la connessione in sicurezza

\end{enumerate}

è considerato non sicuro usare una chiave per piu di una operazione  di cifratura, per cui la key derivation crea 4 chiavi
\begin{itemize}
\item Kc per dati da client a server
\item Mc  chiave MAC per dati da client a server
\item Ks per dati da server a client
\item Ms chiave MAC da server a client
\end{itemize}

step di handshake:

1. client sends list of algorithms it supports, along with
client nonce
2. server chooses algorithms from list; sends back:
choice + certificate + server nonce
3. client verifies certificate, extracts server’s public
key, generates pre master secret, encrypts with
server’s public key, sends to server
4. client and server independently compute encryption
and MAC keys from pre master secret and nonces
5. client sends a MAC of all the handshake messages
6. server sends a MAC of all the handshake messages













\end{document}